% Overlapping subproblems. A naive recursion on the value function expands the
% same state many times. Here we show the recursion for V_t(s) as a tree. Nodes
% with the same (stage, SoC) label, shaded the same color, are the same
% subproblem. A table stores each one so it is solved only once.
\begin{tikzpicture}[
    level distance=10mm,
    sibling distance=15mm,
    every node/.style={font=\scriptsize},
    sub/.style={draw, rounded corners=2pt, inner sep=2pt, minimum size=6mm},
    rep/.style={sub, fill=orange!25},
    repb/.style={sub, fill=cyan!20},
    norm/.style={sub, fill=blue!6},
    edge from parent/.style={draw, black!55, -{Stealth[length=4pt]}},
]
\node[norm] {$V_t(s)$}
  child {node[norm] {$V_{t+1}(s{-}1)$}
    child {node[rep] {$V_{t+2}(s{-}1)$}}
    child {node[repb] {$V_{t+2}(s{-}2)$}}
  }
  child {node[norm] {$V_{t+1}(s)$}
    child {node[rep] {$V_{t+2}(s{-}1)$}}
    child {node[norm] {$V_{t+2}(s)$}}
  }
  child {node[norm] {$V_{t+1}(s{+}1)$}
    child {node[norm] {$V_{t+2}(s)$}}
    child {node[repb] {$V_{t+2}(s{-}2)$}}
  };
% Note about the repeats.
\node[font=\scriptsize, align=left, anchor=north west]
  at (3.0,-0.2) {Same color, same\\subproblem. The DP\\table solves each\\one only once.};
\end{tikzpicture}
